
\section{Centered Bound for the Shifted First-Order Perturbation}

% This chapter is partly preliminary. The last-iterate RR discussion at the end
% is motivational and is not used in the final Berry--Esseen assembly; the
% downstream input retained by Chapter 4 is the centered shifted
% first-order bound proved below. The proof uses the deterministic-product
% perturbation expansion from Samsonov et al. (2025, Proposition 9), specialized
% to the constant-stepsize setting, and isolates the future-centered bilinear
% concentration input used in that argument.

Let
\begin{equation}
\begin{aligned}
B = I - \alpha \overline{A}.
\end{aligned}
\end{equation}
Assume throughout that \(0 < \alpha \le \alpha_\infty\). Then
\begin{equation}
\begin{aligned}
\lVert B^m \rVert_Q \le (1 - \alpha a)^{m / 2},
\quad m \ge 0.
\end{aligned}
\end{equation}
Set also
\begin{equation}
\begin{aligned}
\alpha_{\text{inv}} \coloneqq \frac{1}{2 \lVert  \overline{A}  \rVert}.
\end{aligned}
\end{equation}
If \(w \le \alpha_{\text{inv}}\), the Neumann series gives
\(\lVert  (I - w \overline{A})^{-1}  \rVert \le 2\).
In Richardson--Romberg applications the single-stepsize estimates are applied
at \(w \in \left\{\alpha, 2 \alpha\right\}\), so the local assumptions are
\(2 \alpha \le \alpha_\infty\) and \(2 \alpha \le \alpha_{\text{inv}}\).
Constants denoted by \(C\) may depend on fixed problem parameters, but not on
\(\alpha\), \(n\), \(p\), or \(q\).

The deterministic-product components are defined by
\begin{equation}
\begin{aligned}
J_n^{(0, \alpha)}
  = B J_{n-1}^{(0, \alpha)} - \alpha \varepsilon(Z_n),
\quad
J_0^{(0, \alpha)} = 0,
\end{aligned}
\end{equation}
and
\begin{equation}
\begin{aligned}
J_n^{(1, \alpha)}
  = B J_{n-1}^{(1, \alpha)}
    - \alpha \widetilde{A}(Z_n) J_{n-1}^{(0, \alpha)},
\quad
J_0^{(1, \alpha)} = 0.
\end{aligned}
\end{equation}
Hence
\begin{equation}
\begin{aligned}
J_n^{(0, \alpha)}
  = - \alpha \sum_{k=1}^n B^{n-k} \varepsilon(Z_k),
\end{aligned}
\end{equation}
and
\begin{equation}
\begin{aligned}
J_n^{(1, \alpha)}
&= - \alpha \sum_{j=1}^n
      B^{n-j} \widetilde{A}(Z_j) J_{j-1}^{(0, \alpha)} \\
% &= \alpha^2 \sum_{1 \leq k < j \leq n}
%      B^{n-j}\widetilde{A}(Z_j)B^{j-1-k}\varepsilon(Z_k) \\
&= \alpha^2 \sum_{k=1}^{n-1} \sum_{r=1}^{n-k}
      B^{n-k-r} \widetilde{A}(Z_{k+r}) B^{r-1} \varepsilon(Z_k).
\end{aligned}
\end{equation}
% The sign in the second line is positive because
% \(J_{j-1}^{(0, \alpha)}\) already contains the minus sign.

Use the shifted version
\begin{equation}
\begin{aligned}
S_n
  = \sum_{t=0}^{n-1}
      B^{n-t} \widetilde{A}(Z_{t+1}) J_t^{(0, \alpha)}.
\end{aligned}
\end{equation}
Then
\begin{equation}
\begin{aligned}
S_n
  = \sum_{j=1}^n B^{n-j+1} \widetilde{A}(Z_j) J_{j-1}^{(0, \alpha)}
  = - \frac{1}{\alpha} \, B \, J_n^{(1, \alpha)}.
\end{aligned}
\end{equation}
Thus the corresponding shifted first-order contribution is
\begin{equation}
\begin{aligned}
T_n^{(1, \alpha)} = - \alpha S_n = B \, J_n^{(1, \alpha)},
\end{aligned}
\end{equation}
so passing from \(T_n^{(1, \alpha)}\) back to \(J_n^{(1, \alpha)}\) uses the local
bound on \(B^{-1}\).

\begin{lemma}\label{lem:future-centered-bilinear-input}
  \textbf{(Samsonov et al., 2025,
  Appendix D.2, Proposition 9)}
  Assume UGE 1 and stationarity. Let \(\mathcal{F}_k = \sigma(Z_0, \dots, Z_k)\).
  For \(1 \le k \le n - 1\) and \(1 \le l \le n - k\), let
  \(g_{k,l} : \mathcal{Z} \to \mathbb{R}^d\) be deterministic functions satisfying
  \(\pi(g_{k,l}) = 0\) and \(\lVert g_{k,l} \rVert_\infty \le \beta_{k,l}\). Let
  \(\xi_k\) be \(\mathcal{F}_k\)-measurable vectors with
  \(\sup_k \lVert \xi_k \rVert_\infty \le M_\xi\). Then, for every \(p \ge 2\),
\begin{equation}
\begin{aligned}
  &\lVert  \sum_{k=1}^{n-1} \left\{
      \sum_{l=1}^{n-k} g_{k,l}(Z_{k+l})^\top \xi_k
      - \mathbb{E} \left[
          \sum_{l=1}^{n-k} g_{k,l}(Z_{k+l})^\top \xi_k
          \, | \, \mathcal{F}_k
        \right]
    \right\}  \rVert_{L_p}
  &\quad \le C p^{3 / 2} t_{\text{mix}}^{1 / 2} M_\xi
     \left(\sum_{k=1}^{n-1} \sum_{l=1}^{n-k} \beta_{k,l}^2\right)^{1 / 2}.
\end{aligned}
\end{equation}
% This is the scalar, constant-stepsize form extracted from Samsonov et al.
% (2025, Appendix D.2, Proposition 9). The centering is conditional with
% respect to \(\mathcal{F}_k\).
\end{lemma}

\begin{lemma}\label{lem:last-shifted-first-order}
  \textbf{(Centered shifted first-order bound.)}
  Assume the Markov chain is started from stationarity, that is, the law of
  \(Z_1\) is \(\pi\), and \(\pi(\widetilde{A}) = 0\). For every deterministic direction
  \(u \in \mathbb{R}^d\) and every \(p \ge 2\),
\begin{equation}
\begin{aligned}
  \lVert u^\top (S_n - \mathbb{E}S_n) \rVert_{L_p}
    \le C \lVert u \rVert \, \lVert \varepsilon \rVert_\infty
      (p^{3 / 2} t_{\text{mix}}^{1 / 2} \frac{1}{a}
        + p^{1 / 2} t_{\text{mix}}^{3 / 2} \sqrt{\frac{\alpha}{a}}).
\end{aligned}
\end{equation}
% Consequently,
% \begin{equation}
% \begin{aligned}
% \lVert u^\top (T_n^{(1, \alpha)} - \mathbb{E}T_n^{(1, \alpha)}) \rVert_{L_p}
%   \le C \alpha \lVert u \rVert \, \lVert \varepsilon \rVert_\infty
%     (p^{3 / 2} t_{\text{mix}}^{1 / 2} \frac{1}{a}
%       + p^{1 / 2} t_{\text{mix}}^{3 / 2} \sqrt{\frac{\alpha}{a}}).
% \end{aligned}
% \end{equation}
\end{lemma}

\textit{Proof.} Since \(J_0^{(0, \alpha)} = 0\), the term \(t = 0\) in \(S_n\) vanishes.
Substituting the explicit formula for \(J_t^{(0, \alpha)}\) gives
\begin{equation}
\begin{aligned}
S_n
&= - \alpha \sum_{t=1}^{n-1} \sum_{k=1}^t
    B^{n-t} \widetilde{A}(Z_{t+1}) B^{t-k} \varepsilon(Z_k)
&= - \alpha \sum_{k=1}^{n-1}
    H_{k+1}^{(w)} \varepsilon(Z_k),
\end{aligned}
\end{equation}
where, after the change of summation index \(l = t - k + 1\),
\begin{equation}
\begin{aligned}
H_{k+1}^{(w)}
  = \sum_{l=1}^{n-k}
      B^{n-k-l+1} \widetilde{A}(Z_{k+l}) B^{l-1}.
\end{aligned}
\end{equation}
% The kernel is future-dependent relative to \(\varepsilon(Z_k)\).

Define
\begin{equation}
\begin{aligned}
\mu_k^{(w)} = \mathbb{E}_\pi \left[H_{k+1}^{(w)} \varepsilon(Z_k)\right].
\end{aligned}
\end{equation}
Then
\begin{equation}
\begin{aligned}
S_n - \mathbb{E}S_n
  = - \alpha \sum_{k=1}^{n-1}
      \left\{H_{k+1}^{(w)} \varepsilon(Z_k) - \mu_k^{(w)}\right\}.
\end{aligned}
\end{equation}

Let \(\mathcal{F}_k = \sigma(Z_0, \dots, Z_k)\). The Markov property gives
\begin{equation}
\begin{aligned}
\mathbb{E}[H_{k+1}^{(w)} \varepsilon(Z_k) | \mathcal{F}_k]
  = v_k^{(w, \epsilon)}(Z_k),
\end{aligned}
\end{equation}
where
\begin{equation}
\begin{aligned}
v_k^{(w)}(z)
  = \sum_{l=1}^{n-k}
      B^{n-k-l+1} (\mathsf{Q}^l \widetilde{A})(z) B^{l-1},
\quad
v_k^{(w, \epsilon)}(z)
  = v_k^{(w)}(z) \varepsilon(z).
\end{aligned}
\end{equation}
Here \(\mathsf{Q}\) is the Markov transition kernel, so
\begin{equation}
\begin{aligned}
(\mathsf{Q}^l \widetilde{A})(z)
  = \int \widetilde{A}(u) \, \mathsf{Q}^l (z, \, d u)
  = \mathbb{E} \left[\widetilde{A}(Z_{k+l}) | Z_k = z\right].
\end{aligned}
\end{equation}

Under stationarity,
\(\mu_k^{(w)} = \pi(v_k^{(w, \epsilon)})\). Therefore
\begin{equation}
\begin{aligned}
S_n - \mathbb{E}S_n = - \alpha (U_M + U_R),
\end{aligned}
\end{equation}
with
\begin{equation}
\begin{aligned}
U_M
  = \sum_{k=1}^{n-1}
      \left\{H_{k+1}^{(w)} \varepsilon(Z_k)
       - v_k^{(w, \epsilon)}(Z_k)\right\}
\end{aligned}
\end{equation}
and
\begin{equation}
\begin{aligned}
U_R
  = \sum_{k=1}^{n-1}
      \left\{v_k^{(w, \epsilon)}(Z_k)
       - \pi(v_k^{(w, \epsilon)})\right\}.
\end{aligned}
\end{equation}
% We bound \(U_R\) and \(U_M\) separately.

Fix a deterministic direction \(u\), \(\lVert u \rVert = 1\); the general statement
follows by multiplying by \(\lVert u \rVert\).

For \(U_R\), since \(\pi(\widetilde{A}) = 0\), UGE gives
\begin{equation}
\begin{aligned}
\lVert (\mathsf{Q}^l \widetilde{A})(z) \rVert
  = \left(\lVert  \int \widetilde{A}(y) \left\{\mathsf{Q}^l (z, \, d y) - \pi(d y)\right\}  \rVert\right)
  \le 2 C_A \Delta(\mathsf{Q}^l)
  \le 2 C_A (1 / 4)^{\left\lfloor l / t_{\text{mix}} \right\rfloor}.
\end{aligned}
\end{equation}
Hence
\begin{equation}
\begin{aligned}
\lVert v_k^{(w)} \rVert_\infty
&\le C \sum_{l=1}^{n-k}
    (1 - \alpha a)^{(n-k-l+1) / 2}
    \Delta(\mathsf{Q}^l)
    (1 - \alpha a)^{(l-1) / 2} 
% &\le C(1-\alpha a)^{(n-k)/2}
%     \sum_{l=1}^{\infty}(1/4)^{\left\lfloor l/t_{\mathrm{mix}}\right\rfloor} \\
&\le C t_{\text{mix}} (1 - \alpha a)^{(n-k) / 2}.
\end{aligned}
\end{equation}
Thus
\begin{equation}
\begin{aligned}
\lVert v_k^{(w, \epsilon)} \rVert_\infty
  \le C t_{\text{mix}} \lVert \varepsilon \rVert_\infty
      (1 - \alpha a)^{(n-k) / 2}.
\end{aligned}
\end{equation}
Applying Markov concentration to the centered functions
\(v_k^{(w, \epsilon)} - \pi(v_k^{(w, \epsilon)})\) gives
\begin{equation}
\begin{aligned}
\lVert u^\top U_R \rVert_{L_p}
&\le C p^{1 / 2} t_{\text{mix}}^{1 / 2}
    \left(\sum_{k=1}^{n-1} \lVert v_k^{(w, \epsilon)} \rVert_\infty^2\right)^{1 / 2}
% &\le C p^{1/2}t_{\mathrm{mix}}^{3/2}\lVert\varepsilon\rVert_\infty
%     \left(\sum_{k=1}^{n-1}(1-\alpha a)^{n-k}\right)^{1/2} \\
&\le C p^{1 / 2} t_{\text{mix}}^{3 / 2}
    \lVert \varepsilon \rVert_\infty \frac{1}{\sqrt{\alpha a}}.
\end{aligned}
\end{equation}

For \(U_M\), unfold the projected matrix kernel through
\((H_{k+1}^{(w)})^\top u\):
\begin{equation}
\begin{aligned}
(H_{k+1}^{(w)})^\top u
  = \sum_{l=1}^{n-k} g_{k,l}(Z_{k+l}),
\quad
g_{k,l}(z)
  = (B^{l-1})^\top \, \widetilde{A}(z)^\top \, (B^{n-k-l+1})^\top u.
\end{aligned}
\end{equation}
Each \(g_{k,l}\) is centered under \(\pi\), and
\begin{equation}
\begin{aligned}
\lVert g_{k,l} \rVert_\infty
  \le C (1 - \alpha a)^{(n-k) / 2},
\end{aligned}
\end{equation}
\begin{equation}
\begin{aligned}
\sum_{l=1}^{n-k} \lVert g_{k,l} \rVert_\infty^2
  \le C (n-k)(1 - \alpha a)^{n-k}.
\end{aligned}
\end{equation}
% Unlike \(U_R\), the conditional-centering estimate keeps the factor \(n-k\).

The conditional expectation in the future-centered estimate is exactly
\begin{equation}
\begin{aligned}
\mathbb{E} \left[
  \sum_{l=1}^{n-k} g_{k,l}(Z_{k+l})^\top \varepsilon(Z_k)
  \, | \, \mathcal{F}_k
\right]
  = u^\top v_k^{(w, \epsilon)}(Z_k).
\end{aligned}
\end{equation}
Apply that estimate with
\(\xi_k = \varepsilon(Z_k)\) and
\(\beta_{k,l} = C (1 - \alpha a)^{(n-k) / 2}\):
\begin{equation}
\begin{aligned}
\lVert u^\top U_M \rVert_{L_p}
&\le C p^{3 / 2} t_{\text{mix}}^{1 / 2} \lVert \varepsilon \rVert_\infty
    \left(\sum_{k=1}^{n-1} \sum_{l=1}^{n-k}
      (1 - \alpha a)^{n-k}\right)^{1 / 2}
% &= C p^{3/2}t_{\mathrm{mix}}^{1/2}\lVert\varepsilon\rVert_\infty
%     \left(\sum_{k=1}^{n-1}(n-k)(1-\alpha a)^{n-k}\right)^{1/2} \\
&\le C p^{3 / 2} t_{\text{mix}}^{1 / 2}
    \lVert \varepsilon \rVert_\infty \frac{1}{\alpha a}.
\end{aligned}
\end{equation}
% The additional factor compared with \(U_R\) comes from the conditional
% centering estimate.

% Finally, since \(S_n - \mathbb{E}S_n = -\alpha (U_M + U_R)\),
\begin{equation}
\begin{aligned}
\lVert u^\top (S_n - \mathbb{E}S_n) \rVert_{L_p}
&\le \alpha \, (\lVert u^\top U_M \rVert_{L_p} + \lVert u^\top U_R \rVert_{L_p})
&\le C \lVert \varepsilon \rVert_\infty
  \left(p^{3 / 2} t_{\text{mix}}^{1 / 2} \frac{1}{a}
    + p^{1 / 2} t_{\text{mix}}^{3 / 2} \sqrt{\frac{\alpha}{a}}\right).
\end{aligned}
\end{equation}
% Restoring \(\lVert u \rVert\) and using \(T_n^{(1, \alpha)} = -\alpha S_n\) proves the
% second display. \hfill \(\square\)
